

%
%
%
%
%
%
\chapter{Conclusions}
This thesis investigates the feasibility of mechanizing chapters 1 and 2 of
\emph{Semantics with Applications: An Appetizer} by~\cite{nielson_semantics_2007}
using the Lean 4 interactive theorem prover.
By evaluating the transition from textbook prose to a mechanized proof,
the research identified the specific structural requirements and pedagogical
costs associated with a "verbatim" formalization approach.

\section{Summary of Findings}
The results suggest that while
the logical hierarchy of the textbook by~\cite{nielson_semantics_2007} translates effectively into Lean,
the implementation requires a deliberate layer of "scaffolding"
to bridge the gap between human intuition and machine rigor.

In Natural Semantics, the mapping is highly direct.
However, the introduction of Structural Operational Semantics (SOS)
requires more administrative overhead.
While the mechanization remains faithful to the textbook's use of
configurations ($\gamma$), Lean requires an explicit and exhaustive
management of the transition system.

The "notation gap" identified in the discussion highlights that while
a human reader can fluidly interpret overloaded symbols,
Lean requires a rigorous definition of operational
priority and type-safe configuration constructors to maintain the
verbatim structure of the rules.

We also identified some Pedagogical Shortcuts,
the mechanization surfaced several implicit assumptions that must be
made explicit for Lean:
\begin{itemize}
    \item \textbf{Computational Identity:}
    The "simple string" assumption for variables must be
    supported by DecidableEq to allow for computable state updates.

    \item \textbf{The Myth of the "Trivial" Step:}
    Perhaps the most significant finding is that no step in a
    formal proof can be dismissed as "straightforward" or "analogous."
    Where a textbook might rely on the reader's intuition to bridge
    small logical gaps, Lean requires every intermediate state and every
    application of transitivity to be explicitly constructed and verified.

    \item \textbf{Exhaustive Reasoning:}
    Lean forces the formal closure of "impossible" cases
    (such as the false branch of a verified true condition)
    which are avoided by the use of external side Conditions
    which are evaluated spearately if not implicitly fufilled.
\end{itemize}



\section{Limitations and Future Work}
This thesis focuses on the foundational chapters of \emph{Semantics with Applications: An Appetizer} by \cite{nielson_semantics_2007}.
To further this research, the following areas are proposed:
\begin{itemize}
    \item \textbf{Tactic Scaffolding:}
    Developing custom Lean tactics to automate the "scaffolding"
    required for SOS configurations, potentially allowing students
    to interact with Lean using the same level of abstraction found
    in the textbook.

    \item \textbf{Advanced Language Features:}
    Extending the verbatim mechanization to include variable binding
    and local scoping to determine if the "scaffolding" approach remains
    viable as the language complexity increases.

    \item \textbf{Pedagogical Impact:}
    Investigating whether the requirement to prove "trivial" steps
    helps students develop a deeper understanding of semantic rigor
    or if the technical overhead acts as a barrier to learning.
\end{itemize}

\noindent
In conclusion, while the structures of the foundational chapter in \cite{nielson_semantics_2007} textbook are a
natural fit for Lean's inductive logic, the transition reveals that
"textbook simplicity" often rests on a foundation of implicit assumptions.
Mechanization does not change the logic, but it demands a total accounting
of the steps that humans usually leave unsaid.
